\scanpage{14}
\mspage{5}
% Archive foliation: 13, lower right.
\uline{Ex 1} $M$ catégorie, avec $\underset{\longrightarrow}{\lim}$ et $\underset{\longleftarrow}{\lim}$
quelconques. On pose
\[
\mathcal H(A)=\uline{\operatorname{Hom}}(A,M)
\]
Alors on sait que les foncteurs
\[
f_M^{*}:\underset{\substack{\Vert\\\mathcal H(B)}}{\uline{\operatorname{Hom}}(B,M)}
\longrightarrow
\underset{\substack{\Vert\\\mathcal H(A)}}{\uline{\operatorname{Hom}}(A,M)}
\]
associés aux $f:A\longrightarrow B$ dans $\mathrm{Cat}$, admettent
des adjoints à g. $f_!^M$ et à dr. $f_*^M$,
Ici $\mathrm{Hot}_{\mathcal H}=M$ -- ce cas est
considéré trivial.
\margin{Ici $\left\{\begin{aligned}
\int_A&=\underset{A}{\underset{\longrightarrow}{\lim}}\\
\prod_A&=\underset{A}{\underset{\longleftarrow}{\lim}}
\end{aligned}\right.$}

\uline{Ex 2} Soit $(M,W)$, avec $W\subset\operatorname{Fl}(M)$.
On définit
\[
\unsure{\mathcal W_A|W_A}\subset\operatorname{Fl}\uline{\operatorname{Hom}}(A,M)
\]
\deleted{\unsure{par}} équiv. fibre par fibre, et on
\inserted{pose}
\[
\mathcal H(A)=W_A^{-1}\uline{\operatorname{Hom}}(A,M)
\]
La formation des $f_{\mathcal H}^{*}$ est évidente.
\margin{Attention, \unsure{il n'est pas clair} ici que
$\mathcal H(\coprod A_i)\xrightarrow{\sim}\prod_I\mathcal H(A_i)$,
\unsure{sauf} pour $I$ fini.}
